\documentclass[12pt]{article}
\usepackage[utf8]{inputenc}
\usepackage[T1]{fontenc}
\usepackage[margin=1in]{geometry}
\usepackage{setspace}
\usepackage{xcolor}
\usepackage[colorlinks=true, urlcolor=blue!60!black, linkcolor=blue!60!black]{hyperref}
\setlength{\parindent}{2em}
\setlength{\parskip}{1.25em}
\renewcommand{\baselinestretch}{1.0}
\begin{document}
\begin{center}
{\LARGE \textit{The Slaughter-House Cases}, 83 U.S. 36 (1873)\par}
\vspace{0.8em}
{\large Case Brief\par}
\end{center}
\vspace{1em}

\section{Memory Jogger}

In its first interpretation of the Reconstruction Amendments, the Supreme Court upheld a 25-year Louisiana slaughterhouse monopoly and narrowly confined the Fourteenth Amendment's Privileges or Immunities Clause to rights owing their existence to the federal government rather than ordinary civil or economic rights of state citizenship (\href{https://www.courtlistener.com/opinion/88661/butchers-benevolent-assn-v-crescent-city-live-stock-landing/}{CourtListener}; \href{https://www.law.cornell.edu/supremecourt/text/83/36}{LII}).

\section{Detailed Case Facts}

On March 8, 1869, the Louisiana legislature enacted ``An act to protect the health of the city of New Orleans, to locate the stock-landings and slaughter-houses, and to incorporate the Crescent City Live-Stock Landing and Slaughter-House Company,'' effective June 1, 1869 (\href{https://www.law.cornell.edu/supremecourt/text/83/36}{83 U.S. at 38--43}). Responding to severe public-health concerns caused by slaughterhouse offal and livestock waste entering the Mississippi River upstream of New Orleans's water supply, the statute prohibited landing, keeping, or slaughtering cattle, beeves, calves, sheep, swine, or other animals anywhere within the parishes of Orleans, Jefferson, and St. Bernard---an area of 1,154 square miles inhabited by roughly 200,000 to 300,000 people---except at facilities designated under the act (\href{https://www.law.cornell.edu/supremecourt/text/83/36}{LII}).

The statute chartered seventeen named corporators as the Crescent City Live-Stock Landing and Slaughter-House Company and granted the corporation a 25-year exclusive privilege to operate all stock-landings, stock-yards, and slaughterhouses in the three-parish district (\href{https://www.law.cornell.edu/supremecourt/text/83/36}{LII}). Section 3 required the company to erect a ``grand slaughter-house'' capable of slaughtering 500 animals per day below the city limits by June 1, 1869. Section 5 ordered all competing stock-landings and slaughterhouses closed, while obligating the company---under a \$250 penalty per refusal---to allow any butcher to slaughter animals inside its facilities upon payment of capped statutory fees (\$1.00 per beef, 50 cents per hog or calf, and 30 cents per sheep, goat, or lamb), plus surrender of the head, feet, gore, and entrails (excluding the heart and liver), and subject to inspection by a gubernatorially appointed officer (\href{https://www.law.cornell.edu/supremecourt/text/83/36}{LII}).

Roughly 1,000 butchers and livestock dealers in the three parishes had previously operated independent establishments (\href{https://www.law.cornell.edu/supremecourt/text/83/36}{LII}). Representing the displaced butchers, former U.S. Supreme Court Justice John A. Campbell challenged the monopoly as violating both the Thirteenth Amendment (by imposing an ``involuntary servitude'' analogous to feudal banalities) and Section 1 of the Fourteenth Amendment (by abridging their privileges or immunities of United States citizenship, depriving them of liberty and property without due process of law, and denying them equal protection of the laws) (\href{https://www.oyez.org/cases/1850-1900/83us36}{Oyez}; \href{https://www.law.cornell.edu/supremecourt/text/83/36}{LII}).

\section{Procedural History}

Hundreds of suits were filed in Louisiana state courts by and against the Crescent City Company, including actions brought by the Butchers' Benevolent Association of New Orleans, counter-suits by the company, and an enforcement action brought by the Louisiana Attorney General (\href{https://www.law.cornell.edu/supremecourt/text/83/36}{LII}). The Supreme Court of Louisiana upheld the 1869 statute in favor of the Crescent City Company in all of the test cases (\href{https://www.oyez.org/cases/1850-1900/83us36}{Oyez}).

Five cases reached the U.S. Supreme Court on writs of error under Section 25 of the Judiciary Act in December 1870; after a preliminary supersedeas motion (reported at 10 Wall. 273) and a disputed 1871 settlement agreement, three consolidated cases remained before the Court: \textit{The Butchers' Benevolent Association of New Orleans v. The Crescent City Live-Stock Landing and Slaughter-House Company} (two actions) and \textit{Paul Esteben et al. v. The State of Louisiana, ex rel. S. Belden, Attorney-General} (\href{https://www.law.cornell.edu/supremecourt/text/83/36}{LII}). The Court denied the company's motion to dismiss, first heard oral argument on January 11, 1872 (with Justice Nelson absent), ordered reargument before a full bench on February 3--5, 1873, and issued its decision on April 14, 1873 (\href{https://www.oyez.org/cases/1850-1900/83us36}{Oyez}; \href{https://www.law.cornell.edu/supremecourt/text/83/36}{LII}).

\section{Judicial Votes}

\begin{itemize}
  \item \textbf{Vote:} 5--4, affirming the judgments of the Supreme Court of Louisiana (\href{https://www.oyez.org/cases/1850-1900/83us36}{Oyez}; \href{https://fedsoc.org/case/the-slaughter-house-cases}{Federalist Society}).
  \item \textbf{Majority (5):} Justice Samuel Freeman Miller (author), joined by Justices Nathan Clifford, David Davis, William Strong, and Ward Hunt (\href{https://www.courtlistener.com/opinion/88661/butchers-benevolent-assn-v-crescent-city-live-stock-landing/}{CourtListener}).
  \item \textbf{Dissents (4):}
    \begin{itemize}
      \item Justice Stephen Johnson Field filed the principal dissenting opinion, joined by Chief Justice Salmon P. Chase and Justices Noah Haynes Swayne and Joseph P. Bradley (\href{https://www.law.cornell.edu/supremecourt/text/83/36}{LII}).
      \item Justice Joseph P. Bradley joined Justice Field's dissent and filed a separate dissenting opinion (\href{https://www.law.cornell.edu/supremecourt/text/83/36}{LII}).
      \item Justice Noah Haynes Swayne joined Justices Field's and Bradley's dissents and filed a separate dissenting opinion (\href{https://www.law.cornell.edu/supremecourt/text/83/36}{LII}).
      \item Chief Justice Salmon P. Chase joined Justice Field's dissent without writing separately (\href{https://www.law.cornell.edu/supremecourt/text/83/36}{LII}).
    \end{itemize}
\end{itemize}

\section{Holding}

\textbf{Judgment:} Affirmed. Louisiana's 1869 statute granting the Crescent City Company an exclusive 25-year charter to operate stock-landings and slaughterhouses in the New Orleans area did not violate the Thirteenth or Fourteenth Amendments (\href{https://www.law.cornell.edu/supremecourt/text/83/36}{LII}).

\textbf{Rule Announced:} The Fourteenth Amendment's Citizenship Clause establishes two distinct forms of citizenship---citizenship of the United States and citizenship of a State---and the Privileges or Immunities Clause protects only those privileges and immunities that owe their existence to the federal government, its national character, its Constitution, or its laws, leaving fundamental civil and economic rights (such as pursuing an ordinary trade) within the primary domain of state police power (\href{https://www.law.cornell.edu/constitution-conan/amendment-14/section-1/slaughter-house}{Constitution Annotated}; \href{https://www.law.cornell.edu/supremecourt/text/83/36}{LII}). Further, the Thirteenth Amendment bars personal servitude rather than property or commercial monopolies, and the Fourteenth Amendment's Due Process and Equal Protection Clauses do not invalidate an ordinary state health-and-safety monopoly (\href{https://www.law.cornell.edu/supremecourt/text/83/36}{LII}).

\textbf{Holding vs.\ Dicta:} The \textit{holding} is that Louisiana's slaughterhouse monopoly did not impose involuntary servitude under the Thirteenth Amendment, did not abridge any privilege or immunity of national citizenship, and did not violate the Due Process or Equal Protection Clauses. Two famous passages in Justice Miller's opinion are \textit{dicta}: (1) his illustrative list of federal privileges and immunities (such as traveling to the seat of government, using navigable waters, protection on the high seas, habeas corpus, and assembling to petition Congress), which he expressly offered only ``lest it should be said that no such privileges and immunities are to be found''; and (2) his skeptical remark that ``[w]e doubt very much whether any action of a State not directed by way of discrimination against the negroes as a class, or on account of their race, will ever be held to come within the purview of [the Equal Protection Clause],'' which the Court later repudiated as equal-protection doctrine expanded (\href{https://www.law.cornell.edu/supremecourt/text/83/36}{LII}).

\section{Analysis of Opinions}

\subsection{Majority Opinion (Justice Miller)}

Justice Miller began by characterizing the Louisiana statute as a classic exercise of the state police power over unwholesome trades and public health, emphasizing that the act did not strip butchers of their trade because the company was legally obligated to admit all butchers at regulated fees (\href{https://www.law.cornell.edu/supremecourt/text/83/36}{LII}). Turning to the Reconstruction Amendments, Miller read the Thirteenth, Fourteenth, and Fifteenth Amendments through a ``one pervading purpose'': securing the freedom and civil equality of formerly enslaved African Americans (\href{https://www.law.cornell.edu/constitution-conan/amendment-14/section-1/slaughter-house}{Constitution Annotated}). He quickly rejected the Thirteenth Amendment claim because ``involuntary servitude'' targets personal bondage, not commercial restrictions on land or trade.

Miller's central textual move rested on the two sentences of Section 1 of the Fourteenth Amendment. Because the Citizenship Clause distinguishes between ``citizens of the United States'' and citizens ``of the State wherein they reside,'' while the next sentence forbids states only from abridging ``the privileges or immunities of citizens of the United States,'' Miller concluded that the framers intentionally excluded the privileges and immunities of \textit{state} citizenship from federal protection (\href{https://www.law.cornell.edu/supremecourt/text/83/36}{LII}). Citing \textit{Corfield v. Coryell} (1823) and \textit{Paul v. Virginia} (1869), he classified ordinary civil rights---including acquiring property and pursuing a livelihood---as privileges of state citizenship protected only by state law and Article IV's interstate nondiscrimination rule. Extending federal judicial and congressional oversight to all such civil rights, Miller warned, would ``fetter and degrade the State governments'' and make the Supreme Court a ``perpetual censor upon all legislation of the States'' (\href{https://www.law.cornell.edu/supremecourt/text/83/36}{LII}). Finally, Miller summarily rejected the due process claim as unsupported by settled construction and confined the Equal Protection Clause primarily to racial discrimination against formerly enslaved persons (\href{https://www.law.cornell.edu/supremecourt/text/83/36}{LII}).

\subsection{Dissenting Opinion (Justice Field, joined by Chief Justice Chase and Justices Swayne and Bradley)}

Justice Field separated the statute's genuine sanitary provisions (locating slaughtering downstream and requiring inspection) from its 25-year monopoly, calling the public-health justification for the monopoly a ``shallow'' pretence (\href{https://www.law.cornell.edu/supremecourt/text/83/36}{LII}). While he did not rest his dissent on the Thirteenth Amendment, he noted that the Civil Rights Act of 1866 reflected an understanding that unequal restrictions on ordinary callings functioned as badges of servitude.

On the Fourteenth Amendment, Field argued that national citizenship is now primary and that the Privileges or Immunities Clause protects the fundamental, inalienable rights of free citizens described in \textit{Corfield v. Coryell}, foremost among them ``the right to pursue a lawful employment in a lawful manner, without other restraint than such as equally affects all persons'' (\href{https://www.law.cornell.edu/supremecourt/text/83/36}{LII}). Relying on English common law (\textit{The Case of Monopolies}, 1602, and the Statute of Monopolies, 1624) and state decisions invalidating municipal monopolies, Field contended that the majority's reading rendered the Privileges or Immunities Clause ``a vain and idle enactment, which accomplished nothing,'' because the Supremacy Clause already protected purely federal rights from state interference (\href{https://www.law.cornell.edu/supremecourt/text/83/36}{LII}).

\subsection{Dissenting Opinion (Justice Bradley)}

Joining Field, Justice Bradley went even further in articulating both substantive economic liberty and Bill of Rights incorporation. Bradley argued that national citizenship is primary under the Fourteenth Amendment and carries with it both the fundamental rights of Englishmen (life, liberty, and private property, including the right to choose a lawful calling) and the core individual rights enumerated in the Constitution and its early amendments---specifically naming habeas corpus, jury trial, free exercise of religion, free speech and press, peaceable assembly, security against unreasonable searches and seizures, and due process (\href{https://www.law.cornell.edu/supremecourt/text/83/36}{LII}). Bradley also explicitly concluded that a monopoly barring a large class of citizens from pursuing a lawful occupation violates both the Due Process Clause (``Their right of choice is a portion of their liberty; their occupation is their property'') and the Equal Protection Clause, dismissing the majority's federalism fears as an ``argument from inconvenience'' that cannot override the national will (\href{https://www.law.cornell.edu/supremecourt/text/83/36}{LII}).

\subsection{Dissenting Opinion (Justice Swayne)}

Justice Swayne joined Field and Bradley and framed the Reconstruction Amendments as a ``new Magna Charta'' marking a deliberate structural shift ``at the opposite pole from the first eleven'' amendments, which had restrained only the federal government (\href{https://www.law.cornell.edu/supremecourt/text/83/36}{LII}). Swayne maintained that the text of Section 1 is plain and universal, covering all citizens and persons without racial limitation, and that ``[l]abor is property'' entitled to protection under both the Privileges or Immunities and Due Process Clauses. Answering the majority's fear of novel federal power, Swayne countered that ``the novelty was known and the measure deliberately adopted,'' warning that the majority's narrow construction ``turns, as it were, what was meant for bread into a stone'' (\href{https://www.law.cornell.edu/supremecourt/text/83/36}{LII}).

\section{Examples: Future Applications}

\begin{enumerate}
  \item \textbf{Same-Side (Upholding State Economic Regulation):} A state legislature enacts a licensing regime limiting the number of commercial waste-hauling permits in a metropolitan area to a single regulated private contractor that must serve all residents at state-capped rates. Under \textit{Slaughter-House}, a displaced hauler's Privileges or Immunities Clause challenge fails because pursuing a local trade is a right of state citizenship subject to state police-power regulation rather than a privilege of national citizenship.
  \item \textbf{Same-Side (Denying Privileges or Immunities Protection to Professional Licensing):} A state bars an otherwise qualified state citizen from admission to the state bar based on state-law eligibility criteria. As the Supreme Court held the very next day in \href{https://supreme.justia.com/cases/federal/us/83/130/}{\textit{Bradwell v. Illinois}, 83 U.S. 130 (1873)}, the right to practice a licensed profession in state courts is not a privilege or immunity of national citizenship under \textit{Slaughter-House} (though modern Equal Protection doctrine would independently invalidate sex-based exclusions).
  \item \textbf{Opposite-Side (Protecting a Right of National Citizenship):} A state imposes a special head tax on residents leaving the state to travel to Washington, D.C., or penalizes citizens for traveling to another state to establish bona fide residence and receive equal welfare benefits. Because Miller explicitly recognized the right to travel to the seat of government and the right of a U.S. citizen to become a citizen of any state on equal terms as privileges of national citizenship (\href{https://www.law.cornell.edu/supremecourt/text/83/36}{LII}), such a law abridges the Privileges or Immunities Clause.
  \item \textbf{Opposite-Side (Interference with Federal Navigational or Statutory Rights):} A state prohibits U.S. citizens from operating federally licensed vessels on navigable interstate waterways or obstructs access to federal courts and land offices within the state. Because these rights ``owe their existence to the Federal government, its National character, its Constitution, or its laws'' (\href{https://www.law.cornell.edu/supremecourt/text/83/36}{LII}), they fall squarely within the narrow class of national privileges protected even under \textit{Slaughter-House}.
  \item \textbf{Fence-Sitter (Incorporation of Enumerated Bill of Rights Guarantees via Privileges or Immunities):} A state citizen challenges a state law infringing an enumerated Bill of Rights guarantee---such as the right to keep and bear arms or freedom from excessive fines---specifically under the Privileges or Immunities Clause rather than the Due Process Clause. While \textit{Slaughter-House} long foreclosed Privileges or Immunities incorporation (leading the plurality in \href{https://supreme.justia.com/cases/federal/us/561/742/}{\textit{McDonald v. City of Chicago}, 561 U.S. 742 (2010)}, to retain Due Process incorporation), Miller's own mention of federal constitutional rights like peaceable assembly and habeas corpus has fueled ongoing scholarly and judicial debate---including concurrences by Justices Thomas and Gorsuch---over whether \textit{Slaughter-House} should be revisited or narrowed to allow incorporation through the Privileges or Immunities Clause.
\end{enumerate}

\section{Critique}

\textbf{Originalist and Textualist Perspectives:} Across ideological lines, originalist scholars largely view \textit{Slaughter-House} as one of the Supreme Court's gravest textual errors for effectively erasing the Privileges or Immunities Clause from the Constitution (\href{https://scholarship.law.georgetown.edu/facpub/2147/}{Barnett \& Bernick, Georgetown Law}; \href{https://fedsoc.org/fedsoc-review/protecting-economic-liberty-in-the-federal-courts-theory-precedent-practice}{Federalist Society}). Scholars such as Randy Barnett and Evan Bernick argue that the 39th Congress and the ratifying public understood ``privileges or immunities of citizens of the United States'' to protect both enumerated Bill of Rights guarantees and fundamental unenumerated civil and economic rights (including the right to pursue an honest livelihood), matching the explanation given by Senator Jacob Howard and the Civil Rights Act of 1866 (\href{https://scholarship.law.nd.edu/ndlr/vol95/iss2/1/}{Notre Dame Law Review}). Other originalists, such as Kurt Lash, contend that the Clause's original meaning was limited to enumerated constitutional rights rather than unenumerated common-law rights---a narrower reading than Barnett and Bernick's, yet one that still rejects Miller's near-nullification of the Clause (\href{https://scholarship.law.georgetown.edu/facpub/2147/}{Barnett \& Bernick}). Justice Clarence Thomas has repeatedly urged overruling \textit{Slaughter-House} to anchor Bill of Rights incorporation in the Privileges or Immunities Clause rather than ``substantive due process,'' which he criticizes as an atextual legal fiction (\href{https://supreme.justia.com/cases/federal/us/561/742/}{\textit{McDonald}, 561 U.S. 742}; \href{https://www.law.cornell.edu/supremecourt/text/17-1091}{\textit{Timbs}, 586 U.S. 146}).

\textbf{Progressive and Doctrinal Perspectives:} Progressive and doctrinal scholars offer a more nuanced critique. Many condemn \textit{Slaughter-House} for inaugurating a judicial retreat from Reconstruction that weakened federal civil-rights enforcement and forced later courts to route Bill of Rights incorporation and unenumerated rights awkwardly through the Due Process Clause (\href{https://www.law.cornell.edu/constitution-conan/amendment-14/section-1/slaughter-house}{Constitution Annotated}; \href{https://columbialawreview.org/content/killing-precedent-the-slaughter-house-constitution/}{Columbia Law Review}). At the same time, many progressive commentators defend Miller's \textit{result}---upholding a public-health regulation enacted by Louisiana's biracial Reconstruction legislature---and note the irony that Field's and Bradley's dissents laid the intellectual foundation for \textit{Lochner}-era economic substantive due process, which federal courts later used to strike down labor and welfare legislation (\href{https://fedsoc.org/case/the-slaughter-house-cases}{Federalist Society}; \href{https://columbialawreview.org/content/killing-precedent-the-slaughter-house-constitution/}{Columbia Law Review}). Revisionist historians, including Bryan Wildenthal and Kevin Newsom, have even argued that Miller's opinion did not necessarily reject Bill of Rights incorporation and that subsequent decisions such as \textit{United States v. Cruikshank} (1876) bore greater blame for hollowing out the Clause (\href{https://columbialawreview.org/content/killing-precedent-the-slaughter-house-constitution/}{Columbia Law Review}).

\textbf{Underlying Competing Values:} The case exposes a foundational tension between \textit{antebellum dual federalism} (preserving broad state police power over local health, safety, and economic regulation without making federal courts ``perpetual censors'' of state legislatures) and the \textit{Reconstruction constitutional revolution} (establishing national citizenship as primary and empowering federal courts and Congress to safeguard fundamental individual liberties against state abuse).

\section{Key Quotations}

\begin{quote}
``We repeat, then, in the light of this recapitulation of events, almost too recent to be called history, but which are familiar to us all; and on the most casual examination of the language of these amendments, no one can fail to be impressed with the one pervading purpose found in them all, lying at the foundation of each, and without which none of them would have been even suggested; we mean the freedom of the slave race, the security and firm establishment of that freedom, and the protection of the newly-made freeman and citizen from the oppressions of those who had formerly exercised unlimited dominion over him.''
\end{quote}
\noindent --- Justice Samuel Freeman Miller, for the Court, 83 U.S. at 71 (\href{https://www.law.cornell.edu/supremecourt/text/83/36}{LII}; \href{https://www.courtlistener.com/opinion/88661/butchers-benevolent-assn-v-crescent-city-live-stock-landing/}{CourtListener}).

\begin{quote}
``And still further, such a construction followed by the reversal of the judgments of the Supreme Court of Louisiana in these cases, would constitute this court a perpetual censor upon all legislation of the States, on the civil rights of their own citizens, with authority to nullify such as it did not approve as consistent with those rights, as they existed at the time of the adoption of this amendment.''
\end{quote}
\noindent --- Justice Samuel Freeman Miller, for the Court, 83 U.S. at 78 (\href{https://www.law.cornell.edu/supremecourt/text/83/36}{LII}).

\begin{quote}
``If this inhibition has no reference to privileges and immunities of this character, but only refers, as held by the majority of the court in their opinion, to such privileges and immunities as were before its adoption specially designated in the Constitution or necessarily implied as belonging to citizens of the United States, it was a vain and idle enactment, which accomplished nothing, and most unnecessarily excited Congress and the people on its passage.''
\end{quote}
\noindent --- Justice Stephen Johnson Field, dissenting, 83 U.S. at 96 (\href{https://www.law.cornell.edu/supremecourt/text/83/36}{LII}).

\begin{quote}
``This right to choose one's calling is an essential part of that liberty which it is the object of government to protect; and a calling, when chosen, is a man's property and right. Liberty and property are not protected where these rights are arbitrarily assailed.''
\end{quote}
\noindent --- Justice Joseph P. Bradley, dissenting, 83 U.S. at 116 (\href{https://www.law.cornell.edu/supremecourt/text/83/36}{LII}).

\begin{quote}
``It is objected that the power conferred is novel and large. The answer is that the novelty was known and the measure deliberately adopted. \ldots\ The construction adopted by the majority of my brethren is, in my judgment, much too narrow. It defeats, by a limitation not anticipated, the intent of those by whom the instrument was framed and of those by whom it was adopted. To the extent of that limitation it turns, as it were, what was meant for bread into a stone.''
\end{quote}
\noindent --- Justice Noah Haynes Swayne, dissenting, 83 U.S. at 129 (\href{https://www.law.cornell.edu/supremecourt/text/83/36}{LII}).


\end{document}
